\section{Audit Log}
\label{sec:audit-log}

\subsection{Definisi dan Konsep Audit Log}
\label{subsec:definisi-audit-log}

Audit log didefinisikan sebagai kumpulan \textit{audit record} yang
tersusun secara kronologis (\cite{iso27789_2021}). Definisi ini
menegaskan dua hal sekaligus: pertama, audit log merupakan kumpulan
catatan, bukan catatan tunggal; dan kedua, urutan kronologis adalah
sifat yang melekat, sehingga setiap \textit{record} di dalamnya
terkait dengan waktu kejadiannya dan dapat ditelusuri berdasarkan
rentang waktu tertentu.

ISO 27789:2021 membedakan empat istilah yang berkaitan dalam hierarki
yang berjenjang, sebagaimana ditunjukkan pada
Gambar~\ref{fig:hierarki-audit}. \textit{Audit event} adalah satu
peristiwa pada sistem yang memicu pembangkitan catatan audit
(\cite{iso27789_2021}), misalnya seorang pengguna membaca data rekam
medis. \textit{Audit record} adalah catatan dari satu peristiwa
spesifik tunggal dalam siklus hidup rekam medis elektronik
(\cite{iso27789_2021}). \textit{Audit log} adalah kumpulan
\textit{audit record} yang tersusun secara kronologis
(\cite{iso27789_2021}). Pada level tertinggi, \textit{audit trail}
adalah kumpulan satu atau lebih \textit{audit log} yang secara
keseluruhan membentuk riwayat lengkap atas suatu entitas atau proses
(\cite{iso27789_2021}). Dengan kata lain, satu \textit{audit event}
menghasilkan satu \textit{audit record}; sejumlah \textit{audit record}
yang tersusun membentuk sebuah \textit{audit log}; dan satu atau lebih
\textit{audit log} secara kolektif membentuk sebuah \textit{audit
trail}.

\begin{figure}[htbp]
\centering
% Ganti dengan diagram hierarki jika tersedia
% \includegraphics[width=0.75\textwidth]{chapter-2/hierarki-audit.png}
\caption{Hierarki Konsep Audit Menurut ISO 27789:2021}
\label{fig:hierarki-audit}
\end{figure}

Tugas Akhir ini berfokus pada level \textit{audit log}: pengembangan
mekanisme pembangkitan \textit{audit record} dari berbagai komponen
sistem, pengumpulannya ke dalam sebuah \textit{audit log} yang
terstruktur, serta penyimpanannya secara \textit{tamper-evident}.
Mekanisme ini merupakan prasyarat sebelum \textit{audit trail} yang
utuh dapat terbentuk dari satu atau lebih \textit{audit log} yang
dihasilkan.

Konsep audit log perlu dibedakan secara tegas dari konsep pencatatan
izin akses (\textit{access grant log}). Audit log bersifat komplementer
terhadap kontrol akses, bukan pengganti maupun bagian darinya: kontrol
akses menentukan kebijakan mengenai siapa \textit{berhak} mengakses
data, sedangkan audit log mencatat apa yang \textit{benar-benar
terjadi} pada sistem, termasuk apabila terjadi penyimpangan dari
kebijakan tersebut (\cite{iso27789_2021}). Perbedaan ini penting untuk
digarisbawahi karena suatu sistem dapat saja memiliki kontrol akses
yang ketat dan mencatat metadata pemberian izin secara lengkap, namun
tetap tidak memiliki audit log yang memadai apabila tidak ada catatan
terpisah mengenai peristiwa akses aktual (\textit{access event}) yang
terjadi setelah izin tersebut diberikan. Dengan kata lain, audit log
idealnya mencatat peristiwa (\textit{event}) sebagai representasi dari
tindakan nyata yang terjadi pada sistem, bukan semata representasi dari
kebijakan atau otorisasi yang berlaku.

% -------------------------------------------------------------------
\subsection{Karakteristik Audit Log yang Baik}
\label{subsec:karakteristik-audit-log}

Agar dapat berfungsi sebagai bukti yang dapat diandalkan, audit log
perlu memenuhi sejumlah karakteristik tertentu. Karakteristik pertama
adalah kemampuan mengidentifikasi pengguna secara tidak ambigu
(\textit{unambiguous identification}), yaitu audit log harus
menyediakan data yang cukup untuk mengidentifikasi secara pasti setiap
pengguna maupun sistem eksternal yang telah mengakses atau menerima
data rekam medis (\cite{iso27789_2021}). Karakteristik kedua adalah
integritas (\textit{integrity}) dan kerahasiaan (\textit{confidentiality})
dari audit log itu sendiri, yaitu jaminan bahwa catatan yang telah
dibuat terlindungi dari upaya perusakan (\textit{tampering}), termasuk
perlunya audit log mencatat setiap tindakan yang dilakukan terhadap
sistem audit itu sendiri, seperti kapan sistem audit tidak berfungsi
atau dinonaktifkan (\cite{iso27789_2021}).

Karakteristik lain yang tidak kalah penting adalah kelengkapan
(\textit{completeness}) dan ketersediaan (\textit{availability}).
Kelengkapan berarti audit log harus mencatat seluruh tindakan yang
relevan terhadap data -- termasuk pembuatan, pembacaan, pembaruan,
maupun pengarsipan -- setiap kali tindakan tersebut dilakukan terhadap
data kesehatan pribadi (\cite{iso27789_2021}), sedangkan ketersediaan
berarti sistem audit harus memastikan entri tercatat setiap kali
sistem informasi kesehatan sedang beroperasi (\cite{iso27789_2021}).
Selain keempat karakteristik tersebut, audit log yang baik juga perlu
mendukung kemampuan deteksi terhadap potensi penyalahgunaan oleh pihak
internal (\textit{insider threat}), mengingat pelaku \textit{insider}
pada umumnya memiliki akses sah terhadap sistem sehingga satu-satunya
cara untuk mendeteksi maupun membuktikan penyalahgunaan tersebut adalah
melalui pencatatan aktivitas yang menyeluruh dan tidak dapat
dimanipulasi oleh pelaku itu sendiri (\cite{cert2022insiderthreat}).
Karakteristik-karakteristik tersebut menjadi acuan utama dalam menilai
apakah suatu mekanisme pencatatan pada sistem informasi, termasuk pada
DecMed, telah dapat disebut sebagai audit log yang memadai.

% -------------------------------------------------------------------
\subsection{Komponen Audit Log}
\label{subsec:komponen-audit-log}

Untuk dapat berfungsi secara menyeluruh, suatu sistem audit log perlu
dibangun melalui beberapa komponen fungsional yang saling berurutan.
Kerangka kerja audit untuk rekam medis elektronik membagi persoalan
menjadi penentuan peristiwa apa saja yang perlu dipicu untuk dicatat
(\textit{trigger events}), format dan isi dari setiap catatan yang
dihasilkan, hingga bagaimana catatan tersebut dikelola secara aman
sepanjang siklus hidupnya (\cite{iso27789_2021}). Pembagian yang
serupa juga ditemukan pada kerangka manajemen log komputer secara
umum, yang memandang manajemen log sebagai suatu infrastruktur yang
digunakan untuk menghasilkan (\textit{generate}), mengirimkan
(\textit{transmit}), menyimpan (\textit{store}), serta menganalisis
(\textit{analyze}) data log (\cite{nist80092}). Berdasarkan kedua
kerangka tersebut, komponen audit log pada Tugas Akhir ini
dikelompokkan menjadi empat bagian: penentuan event, pengumpulan
event, penyimpanan event, dan pembacaan event.

\subsubsection{Penentuan Event}
\label{subsubsec:penentuan-event}

Komponen pertama adalah penentuan peristiwa (\textit{event}) apa saja
yang perlu memicu pencatatan pada audit log. Peristiwa pemicu
(\textit{trigger events}) pada sistem rekam medis elektronik pada
umumnya ditentukan berdasarkan skala, tujuan, serta kebijakan privasi
dan keamanan dari masing-masing sistem informasi kesehatan, dengan
cakupan minimal berupa peristiwa akses (\textit{access events}) dan
peristiwa kueri (\textit{query events}) terhadap data kesehatan pribadi
(\cite{iso27789_2021}). Standar tersebut turut memberikan contoh
peristiwa yang secara eksplisit berada di luar cakupan pencatatan
audit rekam medis, seperti peristiwa mulai/berhenti aplikasi,
peristiwa autentikasi pengguna, maupun peristiwa yang dihasilkan oleh
sistem operasi (\cite{iso27789_2021}).

Di luar peristiwa akses dan kueri, penentuan \textit{trigger events}
pada suatu sistem juga dapat diturunkan melalui pendekatan pemodelan
ancaman (\textit{threat modeling}), seperti metode STRIDE, yang
mengidentifikasi ancaman keamanan pada suatu sistem secara sistematis
dan dapat digunakan sebagai dasar penentuan kebutuhan pencatatan
(\cite{laksono2021stride}).

\subsubsection{Konten Wajib dan Konten Tambahan pada Audit Record}
\label{subsubsec:konten-event}

Setiap \textit{audit record}, terlepas dari jenis peristiwa apa yang
dicatatnya, memuat sekumpulan bidang data minimal yang bersifat wajib
(\textit{mandatory}). Format umum suatu \textit{audit record} pada
rekam medis elektronik sekurang-kurangnya memuat identitas unik dari
peristiwa yang terjadi (\textit{EventID}), jenis tindakan yang
dilakukan (\textit{EventActionCode}), waktu terjadinya peristiwa
(\textit{EventDateTime}), identitas pengguna atau proses yang melakukan
tindakan (\textit{UserID}), serta identitas objek data yang menjadi
sasaran tindakan tersebut (\textit{ParticipantObjectID})
(\cite{iso27789_2021}).

Di samping bidang wajib tersebut, \textit{audit record} dapat pula
memuat konten tambahan (\textit{optional}) yang bersifat kontekstual
bergantung pada jenis \textit{event}. Bidang tambahan ini dapat
mencakup informasi mengenai hasil dari tindakan yang dilakukan
(\textit{outcome}), parameter permintaan yang dikirimkan, maupun
atribut bukti digital yang spesifik diperlukan untuk keperluan
investigasi ancaman tertentu (\cite{iso27789_2021}).

\subsubsection{Pengumpulan Event}
\label{subsubsec:pengumpulan-event}

Komponen kedua adalah pengumpulan event, yaitu bagaimana setiap
peristiwa yang terjadi pada sistem direkam menjadi suatu
\textit{audit record} dengan format dan isi yang konsisten.
Pengumpulan event pada praktiknya menghadapi tantangan tersendiri,
terutama pada sistem dengan banyak sumber log (\textit{log sources})
yang berbeda. Setiap sumber log cenderung mencatat informasi dengan
format dan kelengkapan yang berbeda-beda, sehingga menyulitkan proses
menghubungkan (\textit{correlate}) peristiwa yang tercatat oleh satu
sumber dengan peristiwa yang tercatat oleh sumber lainnya
(\cite{nist80092}).

\subsubsection{Penyimpanan Event}
\label{subsubsec:penyimpanan-event}

Komponen ketiga adalah penyimpanan event, yaitu bagaimana catatan
peristiwa yang telah dikumpulkan disimpan dan dikelola secara aman
sepanjang siklus hidupnya. Audit log perlu menyediakan langkah
pengamanan yang cukup untuk melindungi isinya dari perusakan
(\textit{tampering}), termasuk mengamankan akses terhadap
\textit{audit record}, mengamankan akses terhadap perangkat audit itu
sendiri, serta mencatat seluruh tindakan yang dilakukan terhadap audit
log dengan menyertakan waktu, jenis tindakan, dan pelakunya
(\cite{iso27789_2021}). Selain itu, \textit{audit record} juga perlu
tunduk pada kebijakan retensi (\textit{retention policy}) yang
idealnya disesuaikan dengan masa berlaku data rekam medis yang
bersangkutan, mengingat \textit{audit record} dapat dibutuhkan sebagai
bukti dalam jangka waktu yang panjang (\cite{iso27789_2021}).

Pada level infrastruktur, pertimbangan penyimpanan log turut mencakup
keseimbangan antara volume data log yang dihasilkan dengan kapasitas
penyimpanan yang tersedia, baik untuk penyimpanan daring
(\textit{online}) maupun luring (\textit{offline}), serta kebutuhan
keamanan atas data tersebut (\cite{nist80092}).

\subsubsection{Pembacaan Event}
\label{subsubsec:pembacaan-event}

Komponen keempat adalah pembacaan event, yaitu bagaimana catatan
peristiwa yang telah tersimpan dapat diakses kembali dan dianalisis
oleh pihak yang berwenang. Akses terhadap audit log harus dikendalikan
secara ketat dan akses itu sendiri menjadi objek audit
(\textit{itself subject to audit}), dengan akses yang idealnya
dilakukan melalui sistem informasi yang dapat menegakkan kontrol
tersebut, bukan akses langsung terhadap berkas log (\cite{iso27789_2021}).

% -------------------------------------------------------------------
\subsection{Kebutuhan Audit Log}
\label{subsec:kebutuhan-audit-log}

Berdasarkan definisi, karakteristik, dan komponen yang telah diuraikan
pada subbab sebelumnya, dapat diidentifikasi kebutuhan-kebutuhan yang
harus dipenuhi oleh suatu sistem audit log agar dapat berfungsi secara
efektif sebagai bukti digital dalam proses audit maupun investigasi
insiden keamanan. Kebutuhan tersebut dapat dikelompokkan menjadi
kebutuhan fungsional dan kebutuhan keamanan.

\subsubsection{Kebutuhan Fungsional}

Kebutuhan fungsional berkaitan dengan kemampuan yang harus dimiliki
sistem audit log untuk menjalankan fungsinya:

\begin{itemize}
    \item \textbf{Pencatatan event yang komprehensif.} Sistem harus
    mampu mencatat seluruh aktivitas yang relevan dari berbagai
    komponen sistem. Cakupan minimal mencakup peristiwa akses dan
    peristiwa kueri terhadap data sensitif (\cite{iso27789_2021}).

    \item \textbf{Format record yang konsisten.} Setiap
    \textit{audit record} harus memiliki struktur dan atribut yang
    seragam, memuat setidaknya informasi mengenai jenis event, waktu
    kejadian, sumber event, hasil event, serta identitas aktor dan
    objek yang terlibat (\cite{iso27789_2021}).

    \item \textbf{Penyimpanan terpusat dan terdedikasi.} Seluruh
    \textit{audit record} harus tersimpan pada repositori khusus yang
    terpisah dari data operasional sistem, sehingga dapat diakses
    secara independen untuk keperluan audit (\cite{iso27789_2021}).

    \item \textbf{Kemampuan pencarian dan analisis.} Sistem harus
    menyediakan fasilitas untuk mengambil dan menganalisis
    \textit{audit record} berdasarkan berbagai kombinasi atribut
    seperti rentang waktu, jenis event, aktor, maupun objek yang
    terlibat (\cite{iso27789_2021}).

    \item \textbf{Retensi jangka panjang.} \textit{Audit record} harus
    dipertahankan selama periode yang sesuai dengan kebutuhan
    organisasi dan regulasi yang berlaku, mempertimbangkan bahwa
    investigasi dapat dilakukan jauh setelah kejadian
    (\cite{iso27789_2021}).
\end{itemize}

\subsubsection{Kebutuhan Keamanan}

Kebutuhan keamanan berkaitan dengan sifat-sifat yang harus dimiliki
\textit{audit record} agar dapat dipercaya sebagai bukti digital:

\begin{itemize}
    \item \textbf{Integritas.} \textit{Audit record} yang telah
    tersimpan harus terlindungi dari modifikasi yang tidak sah.
    Setiap perubahan pada \textit{audit record} harus dapat terdeteksi
    (\cite{iso27789_2021}).

    \item \textbf{Keaslian (\textit{authenticity}).} Harus tersedia
    mekanisme untuk memverifikasi bahwa \textit{audit record}
    benar-benar dibangkitkan oleh komponen sistem yang diklaim sebagai
    sumbernya, bukan oleh pihak lain yang berpura-pura menjadi sumber
    tersebut (\cite{iso27789_2021}).

    \item \textbf{Non-repudiation.} \textit{Audit record} harus dapat
    digunakan untuk membuktikan bahwa suatu aktivitas benar-benar
    terjadi dan dilakukan oleh entitas tertentu, sehingga entitas
    tersebut tidak dapat menyangkal keterlibatannya
    (\cite{shostack2014}, \cite{iso27789_2021}).

    \item \textbf{Sifat \textit{append-only}.} Mekanisme penyimpanan
    \textit{audit record} harus bersifat tambah-saja, artinya entri
    yang telah tercatat tidak dapat diubah maupun dihapus. Perubahan
    kondisi akibat aktivitas berikutnya dicatat sebagai entri baru
    yang terpisah, bukan sebagai pembaruan entri lama
    (\cite{iso27789_2021}).

    \item \textbf{Akses terkendali.} Akses terhadap \textit{audit
    record} harus dikendalikan secara ketat dan akses itu sendiri
    harus menjadi objek audit, sehingga tidak ada pihak yang dapat
    mengakses maupun memanipulasi \textit{audit record} tanpa
    meninggalkan jejak (\cite{iso27789_2021}).
\end{itemize}

Kebutuhan-kebutuhan tersebut bersifat saling melengkapi: sistem yang
hanya memenuhi kebutuhan fungsional tanpa memenuhi kebutuhan keamanan
tidak dapat diandalkan sebagai bukti digital, sedangkan sistem yang
memenuhi kebutuhan keamanan namun tidak memenuhi kebutuhan fungsional
tidak dapat digunakan untuk investigasi yang komprehensif. Kebutuhan
inilah yang akan digunakan sebagai tolok ukur dalam mengevaluasi
kondisi audit log pada sistem DecMed di Bab III.

% -------------------------------------------------------------------
\subsection{Mekanisme \textit{Tamper-Evident} pada Audit Log}
\label{subsec:mekanisme-tamper-evident}

Untuk memenuhi kebutuhan keamanan yang telah diidentifikasi pada
Subbab~\ref{subsec:kebutuhan-audit-log}, khususnya integritas,
keaslian, dan non-repudiation, terdapat beberapa mekanisme kriptografis
yang lazim digunakan dalam perancangan sistem audit log.

\subsubsection{\textit{Hash-Chaining} untuk Integritas Berurutan}

\textit{Hash-chaining} adalah mekanisme yang menjamin integritas
berurutan (\textit{sequential integrity}) dari sebuah audit log dengan
cara menyertakan nilai \textit{hash} kriptografis dari
\textit{record} sebelumnya ke dalam setiap \textit{record} baru yang
dibuat. Dengan demikian, setiap \textit{audit record} secara kriptografis
bergantung pada seluruh \textit{record} yang mendahuluinya
(\cite{nist80092}).

Implikasi dari mekanisme ini adalah bahwa modifikasi atas satu
\textit{record} dalam rantai -- baik dengan mengubah isi
\textit{record}, menghapusnya, maupun menyisipkan \textit{record}
baru di tengah rantai -- akan menyebabkan nilai \textit{hash} dari
seluruh \textit{record} berikutnya menjadi tidak valid. Ketidaksesuaian
ini dapat dideteksi dengan memverifikasi ulang rantai \textit{hash}
dari awal hingga akhir, sehingga audit log bersifat
\textit{tamper-evident}: perusakan yang terjadi dapat dideteksi meski
tidak dapat dicegah secara langsung pada tingkat berkas lokal.

\subsubsection{Tanda Tangan Digital untuk Keaslian dan Non-Repudiation}

Tanda tangan digital merupakan mekanisme kriptografis yang memungkinkan
suatu entitas membuktikan kepemilikan atas sebuah pesan atau dokumen.
Prinsip kerjanya adalah: pengirim menandatangani data menggunakan
\textit{private key} miliknya, menghasilkan tanda tangan yang dapat
diverifikasi oleh siapa pun yang memiliki \textit{public key}
bersesuaian (\cite{shostack2014}).

Dalam konteks audit log, tanda tangan digital memenuhi dua kebutuhan
sekaligus. \textit{Pertama}, keaslian (\textit{authenticity}): karena
hanya pemegang \textit{private key} yang dapat menghasilkan tanda
tangan yang valid untuk kunci publik tertentu, verifikasi tanda tangan
membuktikan bahwa audit event benar-benar dibangkitkan oleh komponen
yang mengklaim sebagai sumbernya. \textit{Kedua}, non-repudiation:
apabila suatu entitas telah menandatangani sebuah event, entitas
tersebut tidak dapat menyangkal pernah membangkitkan event tersebut
selama \textit{private key}-nya tidak dikompromikan
(\cite{shostack2014}).

Salah satu skema tanda tangan digital yang banyak digunakan untuk
keperluan ini adalah Ed25519, yang berbasis kurva eliptik
Edwards25519. Skema ini bersifat deterministik (tanda tangan yang
dihasilkan selalu sama untuk input yang sama), memiliki ukuran kunci
dan tanda tangan yang kecil (32 byte untuk kunci publik, 64 byte
untuk tanda tangan), dan efisien secara komputasi, menjadikannya
pilihan yang sesuai untuk lingkungan yang membutuhkan verifikasi
tanda tangan dalam volume tinggi.

\subsubsection{Rotasi dan Pengarsipan Log}

Rotasi log adalah praktik manajemen log di mana berkas log aktif
secara berkala diganti dengan berkas baru, sementara berkas lama
diarsipkan untuk keperluan retensi (\cite{nist80092}). Praktik ini
memiliki tiga manfaat utama dalam konteks audit log.

\textit{Pertama}, manajemen ukuran: berkas log yang terus tumbuh
tanpa batas akan menyulitkan pengelolaan dan dapat menghabiskan
kapasitas penyimpanan. Dengan rotasi berkala, setiap berkas arsip
memiliki ukuran yang dapat diprediksi. \textit{Kedua}, efisiensi
analisis: investigasi seringkali berfokus pada rentang waktu
tertentu; berkas arsip yang terorganisasi berdasarkan periode waktu
memudahkan pencarian data yang relevan tanpa harus memindai seluruh
riwayat log. \textit{Ketiga}, keamanan penyimpanan: berkas arsip yang
telah dirotasi dapat segera dipindahkan ke media penyimpanan yang
lebih aman dan tahan lama, terpisah dari lingkungan operasional yang
lebih rentan terhadap gangguan (\cite{nist80092}).

Kombinasi antara \textit{hash-chaining} pada level \textit{record},
tanda tangan digital untuk keaslian \textit{event}, dan rotasi log
untuk pengarsipan berkala membentuk tiga lapisan mekanisme yang saling
melengkapi dalam menjamin integritas, keaslian, dan ketersediaan
jangka panjang dari sebuah sistem audit log yang bersifat
\textit{tamper-evident}.